\chapter{Implementation and Results}
\label{chapter:implementation_and_results}

\section{Development Environment and Technologies}


\section{Implementation of Hexagonal Architecture}



Diagrama de Clases (Arquitectura Hexagonal)

Por qué: Para mostrar cómo el Multi-Format Raster File Input Adapter implementa la interfaz (puerto) necesaria para que el Core pueda leer tanto IDRISI como HDF5.
Qué debe mostrar: La relación de herencia o implementación entre el adaptador y el puerto.

Resumen de Recomendaciones:
1. Herramientas: Usa Draw.io para el Pipeline de Datos (es más visual y libre) y PlantUML para el diagrama de secuencia (es más riguroso para UML).

2. Formatos: Exporta siempre a PDF si usas pdflatex. Esto mantendrá los esquemas vectorizados, permitiendo hacer zoom sin que se vean píxeles, lo cual es vital para la calidad de un TFM.

3. Consistencia: Asegúrate de que los nombres de las clases en el esquema (ej. TemporalLayerManager) coincidan exactamente con los que usas en el texto y en el código.


\subsection{Input}

\subsubsection{Multi-Format Raster File Input Adapter}

Supporting IDRISI and HDF5 formats. (Aquí explicamos que un solo adaptador gestiona ambos).

\subsection{Output}

\subsection{State Updater}

\section{Implementation of the Simulation Core}

\subsection{Algorithm for updating states using tensor operations}

\subsection{Performance optimization and parallelism}

\subsection{Tensor-based Grid Structure}

\subsection{The TemporalLayerManager: Efficient Memory Handling for Dynamic Data}

(Aquí documentamos la lógica de actualización por pasos)

Diagrama de Secuencia UML (Lógica de Ejecución)

Qué debe mostrar:

- La línea de vida del Core (que avanza en el tiempo).

- La línea de vida del TemporalLayerManager (que valida el tiempo).

- La llamada al InputAdapter para leer el "hyperslab" del HDF5.

- La actualización del Tensor en el Grid.

\begin{figure}[htbp]
    \centering
    \includegraphics[width=0.8\textwidth]{figures/sequence_rainfall_update.pdf}
    \caption{UML Sequence Diagram: Rainfall tensor update via TemporalLayerManager.}
    \label{fig:seq_update}
\end{figure}

\paragraph{Logic Workflow:} 
\begin{enumerate}
    \item \textbf{Subscription}: During initialization, the rainfall layer is registered in the 
    TemporalLayerManager with a specific update frequency (e.g., every 600 simulation seconds).
    \item \textbf{Step Validation}: In each iteration of the Core, the manager checks if the 
    current step requires a data refresh.
    \item \textbf{Update Request}: If an update is needed, the manager calls the update method.
    \item \textbf{HDF5 Access}: The manager uses the HighFive functions to perform a Hyperslab selection.
    \item \textbf{Tensor Update}: The buffer is moved into the specific channel of the Core's grid tensor.
\end{enumerate}

\section{Data Extraction and Ingestion Process (Data Pipeline)}

\subsection{Elevation Layer}

\subsection{Water Depth Layer}

\subsection{Roughness Layer}

\subsection{Rainfall Layer}

The implementation of the rainfall layer follows a decoupled pipeline where data acquisition and spatial 
normalization are performed offline using Python, while the real-time ingestion is managed by the C++ Core.

\begin{figure}[htbp]
    \centering
    % Esquema del Pipeline de Datos (Diagrama de Flujo)
    % Aquí iría tu export de Draw.io en PDF o PNG
    % Qué debe mostrar:
    %
    %    - Las fuentes externas (APIs de sensores).
    %
    %    - Los scripts de Python como bloques de procesamiento.
    %
    %    - El archivo HDF5 como el "almacén" final que consume el simulador.
    \includegraphics[width=0.9\textwidth]{figures/rainfall_pipeline.pdf} 
    \caption{Rainfall Data Pipeline: From physical sensors to HDF5 serialization.}
    \label{fig:rainfall_pipeline}
\end{figure}

\paragraph{A. Pre-processing Pipeline (Python ETL)} 

Before the simulation begins, a series of Python scripts perform the \textbf{Extract, Transform, and Load (ETL)} 
process. This approach is chosen to leverage Python's extensive geospatial libraries (such as Pandas, 
Xarray, and Scipy).

\begin{enumerate}
    \item \textbf{Extraction}: The scripts interface with the API or data repositories of the Hydrographic 
    Confederations to retrieve raw meteorological time-series from physical sensors.
    \item \textbf{Transformation (Spatial Interpolation)}: Since rainfall sensors are point-based, the scripts 
    apply a spatial interpolation algorithm (e.g., Inverse Distance Weighting). To ensure consistency, the 
    simulation's DEM is used as a spatial mask; the interpolation is clipped and resampled to match the exact 
    grid dimensions and resolution of the elevation layer.
    \item \textbf{Loading (HDF5)}: The resulting temporal sequence of rasters is serialized into a single 
    HDF5 file. The data is organized into a 3D dataset (time, x, y), optimized for chunked access.
\end{enumerate}

\noindent The following table details the scripts involved in the ETL (Extract, Transform, Load) process:

\begin{table}[H]
\centering
\begin{tabular}{|c|c|c|}
\hline
\textbf{Script} & \textbf{Libraries} & \textbf{Core Responsibility} \\ \hline
\texttt{rain\_fetcher.py}      & \texttt{requests, pandas}              & Automated retrieval of sensor time-series from Hydrographic Confederation APIs.          \\ \hline
\texttt{spatial\_aligner.py}    & \texttt{gdal, scipy}              & Spatial interpolation of point-data and clipping to the DEM's extent/resolution.          \\ \hline
\texttt{hdf5\_optimizer.py}      & \texttt{h5py, numpy}              & Serialization of the raster sequence into a chunked 3D HDF5 dataset ($T, X, Y$).         \\ \hline
\end{tabular}
\caption{Python scripts for Rainfall data pre-processing}
\label{tab:python-scripts-rainfall-data-pre-processing}
\end{table}

\paragraph{B. Core Ingestion and Memory Management} 

In the implementation phase, the TemporalLayerManager class acts as the orchestrator for all time-varying 
data. This component ensures that the Core only allocates memory for the rainfall data that is currently 
relevant to the simulation time $t$.

\section{Testing and Validation}

\subsection{Test scenario: Simulation of a DANA event}

\subsection{Analysis of results and comparison with historical data or theoretical models}



Pruebas, validación y los resultados de la simulación.
